% foundations/information_modelling.tex

\chapter{Information Modelling}

\label{chap:information-modelling}

% ============================================================
\section{Information Modelling}
% ============================================================

Engineering projects rely on information models to describe, organize,
exchange, and interpret engineering knowledge.

An information model is therefore not merely a data structure. It
establishes the conceptual framework within which engineering knowledge
is represented, related, and interpreted.

Different engineering disciplines organize information around different
primary concepts. Buildings, manufacturing systems, transportation
networks, and operational systems each possess characteristic structures
that naturally determine how engineering information is organized.

Alignment Information Modelling (AIM) follows this principle. Rather
than adopting physical assets as its primary organizing concept, AIM
places the railway alignment at the centre of the engineering model.

This perspective reflects the observation that railway systems are
fundamentally governed by their intrinsic geometry. Assets, operational
states, engineering observations, and realization processes are
understood relative to the alignment rather than independently of it.

% ============================================================
\section{Information Modelling Paradigms}
% ============================================================

\begin{definition}[Process Information Modelling]
\label{def:process-information-modelling}

Process Information Modelling (PIM) denotes the general class of
engineering information modelling approaches that organize engineering
knowledge around engineering processes, state evolution, operational
behaviour, and lifecycle relationships rather than exclusively around
physical assets.

Within this thesis, PIM is used as a conceptual umbrella rather than as
an external standard.

\end{definition}

\begin{definition}[Building Information Modelling]
\label{def:building-information-modelling}

Building Information Modelling (BIM) denotes an asset-oriented
information modelling paradigm in which physical objects, their
properties, relationships, and lifecycle information form the primary
organizing structure.

\end{definition}

\begin{definition}[Alignment Information Modelling]
\label{def:alignment-based-information-modelling}

Alignment Information Modelling (AIM) denotes an alignment-oriented
information modelling paradigm in which the railway alignment provides
the primary engineering reference for organizing geometric, spatial,
operational, and engineering information.

\end{definition}

PIM, BIM, and AIM are therefore not competing concepts. They represent
different modelling paradigms emphasizing different organizing
principles appropriate for different engineering domains.

% ============================================================
\section{Why Alignment Information Modelling?}
% ============================================================

Every information model implicitly answers a fundamental engineering
question:

\begin{quote}
What is the primary concept around which engineering knowledge is
organized?
\end{quote}

For buildings, the natural answer is the physical asset. Consequently,
BIM organizes engineering knowledge around objects, assemblies,
components, and their relationships.

Railway systems exhibit a fundamentally different dependency structure.
Vehicle motion, operational behaviour, cant, speed, visibility,
admissible states, and many engineering constraints are determined
primarily by the geometry of the alignment rather than by individual
assets.

The alignment therefore represents more than another engineering object.
It establishes the geometric framework within which railway engineering
takes place.

For this reason, AIM adopts the alignment as the primary organizing
concept of the engineering information model.

\begin{definition}[Primary modelling entity]
\label{def:primary-modelling-entity}

Within Alignment Information Modelling, the railway alignment is the
primary modelling entity.

Engineering information is interpreted relative to the alignment, its
intrinsic geometry, its reference systems, and the engineering states
derived from them.

\end{definition}

This decision does not diminish the importance of engineering assets.
Bridges, tunnels, switches, signals, stations, and other infrastructure
objects remain essential engineering objects.

Their interpretation, however, is no longer independent of the
alignment. They are understood through their relationship to the
alignment and the engineering context established by it.

% ============================================================
\section{Role within this Thesis}
% ============================================================

The preceding chapters established the engineering motivation and
introduced the conceptual vocabulary provided by the Knowledge Kernel.

The purpose of the present chapter is to position Alignment Information
Modelling as the engineering framework within which these concepts are
applied.

The mathematical foundations developed in the following chapters do not
replace these engineering concepts. They provide the formal language
required to describe them rigorously.

Subsequent parts of the thesis progressively develop intrinsic geometry,
metric realization, spatial interpretation, and engineering realization
upon this conceptual foundation.

% ============================================================
\section{Summary}
% ============================================================

\begin{summary}

Engineering information models differ primarily in the concepts around
which engineering knowledge is organized.

BIM adopts physical assets as its primary organizing principle.

AIM adopts the railway alignment as the primary organizing principle.

This alignment-centred perspective provides the conceptual foundation
for the mathematical development that follows and, ultimately, for the
engineering realization of railway systems presented throughout the
remainder of the thesis.

\end{summary}
